\documentclass[tikz,border=5pt]{standalone}
\usepackage[utf8]{inputenc}
\usepackage[T1]{fontenc}
\usepackage[german]{babel}
\usepackage{amsmath, amssymb}
\usepackage{siunitx}
\usepackage{tikz}
\usepackage{pgfplots}
\usepackage{xcolor}
\pgfplotsset{compat=1.18}

\begin{document}
\begin{tikzpicture}
	\begin{axis}[
		width=13cm, height=7cm,
		xlabel={Phase $\omega t$},
		ylabel={Normierte Amplitude},
		grid=both,
		grid style={line width=0.3pt, draw=gray!30},
		xtick={0, 1.5708, 3.1416, 4.7124, 6.2832},
		xticklabels={$0$,$\frac{\pi}{2}$,$\pi$,$\frac{3\pi}{2}$,$2\pi$},
		legend pos=north east,
		xmin=0, xmax=6.2832
	]
		% Spule: u_L führt, i_L folgt
		\addplot[color=red!75!black, very thick, domain=0:6.2832, samples=300]
			{cos(deg(x))};
		\addlegendentry{$u_L(t) = U_0\cos(\omega t)$ — Spannung Spule}
		\addplot[color=blue!75!black, thick, domain=0:6.2832, samples=300]
			{sin(deg(x))};
		\addlegendentry{$i_L(t) = I_0\sin(\omega t)$ — Strom Spule \textbf{(eilt nach)}}
		% Kondensator: i_C führt, u_C folgt
		\addplot[color=orange!85!black, dashed, very thick, domain=0:6.2832, samples=300]
			{sin(deg(x))};
		\addlegendentry{$u_C(t) = U_0\sin(\omega t)$ — Spannung Kondensator}
		\addplot[color=teal!75!black, dashed, thick, domain=0:6.2832, samples=300]
			{cos(deg(x))};
		\addlegendentry{$i_C(t) = I_0\cos(\omega t)$ — Strom Kondensator \textbf{(eilt nach — kausale Wirkung)}}
		% Markierung Viertelperiode
		\draw[dotted, gray!70, thick] (axis cs:1.5708,-1.1) -- (axis cs:1.5708,1.1)
			node[above, font=\small, color=gray!60]{\small $T/4$};
	\end{axis}
\end{tikzpicture}
\end{document}
